\documentclass{beamer}
\usepackage{xeCJK}
\setmainfont{WenQuanYi Micro Hei}
\begin{document}

\section{QuickstartÂ}

\begin{frame}[t]{QuickstartÂ}

Eager to get started?  This page gives a good introduction to Flask.  It
assumes you already have Flask installed.  If you do not, head over to the
Installation section.

\end{frame}

\subsection{A Minimal ApplicationÂ}

\begin{frame}[t]{A Minimal ApplicationÂ}

A minimal Flask application looks something like this:

\end{frame}


\begin{frame}[t,fragile]{A Minimal ApplicationÂ}


\begin{verbatim}

from flask import Flask
app = Flask(__name__)

@app.route('/')
def hello_world():
    return 'Hello, World!'


\end{verbatim}


\end{frame}


\begin{frame}[t]{A Minimal ApplicationÂ}

So what did that code do?

\end{frame}


\begin{frame}[t]{A Minimal ApplicationÂ}

Just save it as hello.py or something similar. Make sure to not call
your application flask.py because this would conflict with Flask
itself.

\end{frame}


\begin{frame}[t]{A Minimal ApplicationÂ}

To run the application you can either use the flask command or
python’s -m switch with Flask.  Before you can do that you need
to tell your terminal the application to work with by exporting the
FLASK\_APP environment variable:

\end{frame}


\begin{frame}[t,fragile]{A Minimal ApplicationÂ}


\begin{verbatim}

$ export FLASK_APP=hello.py
$ flask run
 * Running on http://127.0.0.1:5000/


\end{verbatim}


\end{frame}


\begin{frame}[t]{A Minimal ApplicationÂ}

If you are on Windows you need to use set instead of export.

\end{frame}


\begin{frame}[t]{A Minimal ApplicationÂ}

Alternatively you can use python -m flask:

\end{frame}


\begin{frame}[t,fragile]{A Minimal ApplicationÂ}


\begin{verbatim}

$ export FLASK_APP=hello.py
$ python -m flask run
 * Running on http://127.0.0.1:5000/


\end{verbatim}


\end{frame}


\begin{frame}[t]{A Minimal ApplicationÂ}

This launches a very simple builtin server, which is good enough for testing
but probably not what you want to use in production. For deployment options see
Deployment Options.

\end{frame}


\begin{frame}[t]{A Minimal ApplicationÂ}

Now head over to http://127.0.0.1:5000/, and you
should see your hello world greeting.

\end{frame}


\begin{frame}[t]{A Minimal ApplicationÂ}

Externally Visible Server

\end{frame}


\begin{frame}[t]{A Minimal ApplicationÂ}

If you run the server you will notice that the server is only accessible
from your own computer, not from any other in the network.  This is the
default because in debugging mode a user of the application can execute
arbitrary Python code on your computer.

\end{frame}


\begin{frame}[t]{A Minimal ApplicationÂ}

If you have the debugger disabled or trust the users on your network,
you can make the server publicly available simply by adding
--host=0.0.0.0 to the command line:

\end{frame}


\begin{frame}[t,fragile]{A Minimal ApplicationÂ}


\begin{verbatim}

flask run --host=0.0.0.0


\end{verbatim}


\end{frame}


\begin{frame}[t]{A Minimal ApplicationÂ}

This tells your operating system to listen on all public IPs.

\end{frame}

\subsection{What to do if the Server does not StartÂ}

\begin{frame}[t]{What to do if the Server does not StartÂ}

In case the python -m flask fails or flask does not exist,
there are multiple reasons this might be the case.  First of all you need
to look at the error message.

\end{frame}


\begin{frame}[t]{What to do if the Server does not StartÂ}

Versions of Flask older than 0.11 use to have different ways to start the
application.  In short, the flask command did not exist, and
neither did python -m flask.  In that case you have two options:
either upgrade to newer Flask versions or have a look at the Development Server
docs to see the alternative method for running a server.

\end{frame}


\begin{frame}[t]{What to do if the Server does not StartÂ}

The FLASK\_APP environment variable is the name of the module to import at
flask run. In case that module is incorrectly named you will get an
import error upon start (or if debug is enabled when you navigate to the
application). It will tell you what it tried to import and why it failed.

\end{frame}


\begin{frame}[t]{What to do if the Server does not StartÂ}

The most common reason is a typo or because you did not actually create an
app object.

\end{frame}

\subsection{Debug ModeÂ}

\begin{frame}[t]{Debug ModeÂ}

(Want to just log errors and stack traces? See Application Errors)

\end{frame}


\begin{frame}[t]{Debug ModeÂ}

The flask script is nice to start a local development server, but
you would have to restart it manually after each change to your code.
That is not very nice and Flask can do better.  If you enable debug
support the server will reload itself on code changes, and it will also
provide you with a helpful debugger if things go wrong.

\end{frame}


\begin{frame}[t]{Debug ModeÂ}

To enable debug mode you can export the FLASK\_DEBUG environment variable
before running the server:

\end{frame}


\begin{frame}[t,fragile]{Debug ModeÂ}


\begin{verbatim}

$ export FLASK_DEBUG=1
$ flask run


\end{verbatim}


\end{frame}


\begin{frame}[t]{Debug ModeÂ}

(On Windows you need to use set instead of export).

\end{frame}


\begin{frame}[t]{Debug ModeÂ}

This does the following things:

\end{frame}


\begin{frame}[t]{Debug ModeÂ}

There are more parameters that are explained in the Development Server docs.

\end{frame}


\begin{frame}[t]{Debug ModeÂ}

Attention

\end{frame}


\begin{frame}[t]{Debug ModeÂ}

Even though the interactive debugger does not work in forking environments
(which makes it nearly impossible to use on production servers), it still
allows the execution of arbitrary code. This makes it a major security risk
and therefore it must never be used on production machines.

\end{frame}


\begin{frame}[t]{Debug ModeÂ}

Screenshot of the debugger in action:

\end{frame}


\begin{frame}[t]{Debug ModeÂ}

Have another debugger in mind? See Working with Debuggers.

\end{frame}

\subsection{RoutingÂ}

\begin{frame}[t]{RoutingÂ}

Modern web applications have beautiful URLs.  This helps people remember
the URLs, which is especially handy for applications that are used from
mobile devices with slower network connections.  If the user can directly
go to the desired page without having to hit the index page it is more
likely they will like the page and come back next time.

\end{frame}


\begin{frame}[t]{RoutingÂ}

As you have seen above, the route() decorator is used to
bind a function to a URL.  Here are some basic examples:

\end{frame}


\begin{frame}[t,fragile]{RoutingÂ}


\begin{verbatim}

@app.route('/')
def index():
    return 'Index Page'

@app.route('/hello')
def hello():
    return 'Hello, World'


\end{verbatim}


\end{frame}


\begin{frame}[t]{RoutingÂ}

But there is more to it!  You can make certain parts of the URL dynamic and
attach multiple rules to a function.

\end{frame}


\begin{frame}[t]{RoutingÂ}

To add variable parts to a URL you can mark these special sections as
<variable\_name>.  Such a part is then passed as a keyword argument to your
function.  Optionally a converter can be used by specifying a rule with
<converter:variable\_name>.  Here are some nice examples:

\end{frame}


\begin{frame}[t,fragile]{RoutingÂ}


\begin{verbatim}

@app.route('/user/<username>')
def show_user_profile(username):
    # show the user profile for that user
    return 'User %s' % username

@app.route('/post/<int:post_id>')
def show_post(post_id):
    # show the post with the given id, the id is an integer
    return 'Post %d' % post_id


\end{verbatim}


\end{frame}


\begin{frame}[t]{RoutingÂ}

The following converters exist:

\end{frame}


\begin{frame}[t]{RoutingÂ}

Unique URLs / Redirection Behavior

\end{frame}


\begin{frame}[t]{RoutingÂ}

Flask’s URL rules are based on Werkzeug’s routing module.  The idea
behind that module is to ensure beautiful and unique URLs based on
precedents laid down by Apache and earlier HTTP servers.

\end{frame}


\begin{frame}[t]{RoutingÂ}

Take these two rules:

\end{frame}


\begin{frame}[t,fragile]{RoutingÂ}


\begin{verbatim}

@app.route('/projects/')
def projects():
    return 'The project page'

@app.route('/about')
def about():
    return 'The about page'


\end{verbatim}


\end{frame}


\begin{frame}[t]{RoutingÂ}

Though they look rather similar, they differ in their use of the trailing
slash in the URL definition.  In the first case, the canonical URL for the
projects endpoint has a trailing slash.  In that sense, it is similar to
a folder on a filesystem.  Accessing it without a trailing slash will cause
Flask to redirect to the canonical URL with the trailing slash.

\end{frame}


\begin{frame}[t]{RoutingÂ}

In the second case, however, the URL is defined without a trailing slash,
rather like the pathname of a file on UNIX-like systems. Accessing the URL
with a trailing slash will produce a 404 “Not Found” error.

\end{frame}


\begin{frame}[t]{RoutingÂ}

This behavior allows relative URLs to continue working even if the trailing
slash is omitted, consistent with how Apache and other servers work.  Also,
the URLs will stay unique, which helps search engines avoid indexing the
same page twice.

\end{frame}


\begin{frame}[t]{RoutingÂ}

If it can match URLs, can Flask also generate them?  Of course it can.  To
build a URL to a specific function you can use the url\_for()
function.  It accepts the name of the function as first argument and a number
of keyword arguments, each corresponding to the variable part of the URL rule.
Unknown variable parts are appended to the URL as query parameters.  Here are
some examples:

\end{frame}


\begin{frame}[t,fragile]{RoutingÂ}


\begin{verbatim}

>>> from flask import Flask, url_for
>>> app = Flask(__name__)
>>> @app.route('/')
... def index(): pass
...
>>> @app.route('/login')
... def login(): pass
...
>>> @app.route('/user/<username>')
... def profile(username): pass
...
>>> with app.test_request_context():
...  print url_for('index')
...  print url_for('login')
...  print url_for('login', next='/')
...  print url_for('profile', username='John Doe')
...
/
/login
/login?next=/
/user/John%20Doe


\end{verbatim}


\end{frame}


\begin{frame}[t]{RoutingÂ}

(This also uses the test\_request\_context() method, explained
below.  It tells Flask to behave as though it is handling a request, even
though we are interacting with it through a Python shell.  Have a look at the
explanation below. Context Locals).

\end{frame}


\begin{frame}[t]{RoutingÂ}

Why would you want to build URLs using the URL reversing function
url\_for() instead of hard-coding them into your templates?
There are three good reasons for this:

\end{frame}


\begin{frame}[t]{RoutingÂ}

HTTP (the protocol web applications are speaking) knows different methods for
accessing URLs.  By default, a route only answers to GET requests, but that
can be changed by providing the methods argument to the
route() decorator.  Here are some examples:

\end{frame}


\begin{frame}[t,fragile]{RoutingÂ}


\begin{verbatim}

from flask import request

@app.route('/login', methods=['GET', 'POST'])
def login():
    if request.method == 'POST':
        do_the_login()
    else:
        show_the_login_form()


\end{verbatim}


\end{frame}


\begin{frame}[t]{RoutingÂ}

If GET is present, HEAD will be added automatically for you.  You
don’t have to deal with that.  It will also make sure that HEAD requests
are handled as the HTTP RFC (the document describing the HTTP
protocol) demands, so you can completely ignore that part of the HTTP
specification.  Likewise, as of Flask 0.6, OPTIONS is implemented for you
automatically as well.

\end{frame}


\begin{frame}[t]{RoutingÂ}

You have no idea what an HTTP method is?  Worry not, here is a quick
introduction to HTTP methods and why they matter:

\end{frame}


\begin{frame}[t]{RoutingÂ}

The HTTP method (also often called “the verb”) tells the server what the
client wants to do with the requested page.  The following methods are
very common:

\end{frame}


\begin{frame}[t]{RoutingÂ}

Now the interesting part is that in HTML4 and XHTML1, the only methods a
form can submit to the server are GET and POST.  But with JavaScript
and future HTML standards you can use the other methods as well.  Furthermore
HTTP has become quite popular lately and browsers are no longer the only
clients that are using HTTP. For instance, many revision control systems
use it.

\end{frame}

\subsection{Static FilesÂ}

\begin{frame}[t]{Static FilesÂ}

Dynamic web applications also need static files.  That’s usually where
the CSS and JavaScript files are coming from.  Ideally your web server is
configured to serve them for you, but during development Flask can do that
as well.  Just create a folder called static in your package or next to
your module and it will be available at /static on the application.

\end{frame}


\begin{frame}[t]{Static FilesÂ}

To generate URLs for static files, use the special 'static' endpoint name:

\end{frame}


\begin{frame}[t,fragile]{Static FilesÂ}


\begin{verbatim}

url_for('static', filename='style.css')


\end{verbatim}


\end{frame}


\begin{frame}[t]{Static FilesÂ}

The file has to be stored on the filesystem as static/style.css.

\end{frame}

\subsection{Rendering TemplatesÂ}

\begin{frame}[t]{Rendering TemplatesÂ}

Generating HTML from within Python is not fun, and actually pretty
cumbersome because you have to do the HTML escaping on your own to keep
the application secure.  Because of that Flask configures the Jinja2 template engine for you automatically.

\end{frame}


\begin{frame}[t]{Rendering TemplatesÂ}

To render a template you can use the render\_template()
method.  All you have to do is provide the name of the template and the
variables you want to pass to the template engine as keyword arguments.
Here’s a simple example of how to render a template:

\end{frame}


\begin{frame}[t,fragile]{Rendering TemplatesÂ}


\begin{verbatim}

from flask import render_template

@app.route('/hello/')
@app.route('/hello/<name>')
def hello(name=None):
    return render_template('hello.html', name=name)


\end{verbatim}


\end{frame}


\begin{frame}[t]{Rendering TemplatesÂ}

Flask will look for templates in the templates folder.  So if your
application is a module, this folder is next to that module, if it’s a
package it’s actually inside your package:

\end{frame}


\begin{frame}[t]{Rendering TemplatesÂ}

Case 1: a module:

\end{frame}


\begin{frame}[t,fragile]{Rendering TemplatesÂ}


\begin{verbatim}

/application.py
/templates
    /hello.html


\end{verbatim}


\end{frame}


\begin{frame}[t]{Rendering TemplatesÂ}

Case 2: a package:

\end{frame}


\begin{frame}[t,fragile]{Rendering TemplatesÂ}


\begin{verbatim}

/application
    /__init__.py
    /templates
        /hello.html


\end{verbatim}


\end{frame}


\begin{frame}[t]{Rendering TemplatesÂ}

For templates you can use the full power of Jinja2 templates.  Head over
to the official Jinja2 Template Documentation for more information.

\end{frame}


\begin{frame}[t]{Rendering TemplatesÂ}

Here is an example template:

\end{frame}


\begin{frame}[t,fragile]{Rendering TemplatesÂ}


\begin{verbatim}

<!doctype html>
<title>Hello from Flask</title>
{% if name %}
  <h1>Hello {{ name }}!</h1>
{% else %}
  <h1>Hello, World!</h1>
{% endif %}


\end{verbatim}


\end{frame}


\begin{frame}[t]{Rendering TemplatesÂ}

Inside templates you also have access to the request,
session and g [1] objects
as well as the get\_flashed\_messages() function.

\end{frame}


\begin{frame}[t]{Rendering TemplatesÂ}

Templates are especially useful if inheritance is used.  If you want to
know how that works, head over to the Template Inheritance pattern
documentation.  Basically template inheritance makes it possible to keep
certain elements on each page (like header, navigation and footer).

\end{frame}


\begin{frame}[t]{Rendering TemplatesÂ}

Automatic escaping is enabled, so if name contains HTML it will be escaped
automatically.  If you can trust a variable and you know that it will be
safe HTML (for example because it came from a module that converts wiki
markup to HTML) you can mark it as safe by using the
Markup class or by using the |safe filter in the
template.  Head over to the Jinja 2 documentation for more examples.

\end{frame}


\begin{frame}[t]{Rendering TemplatesÂ}

Here is a basic introduction to how the Markup class works:

\end{frame}


\begin{frame}[t,fragile]{Rendering TemplatesÂ}


\begin{verbatim}

>>> from flask import Markup
>>> Markup('<strong>Hello %s!</strong>') % '<blink>hacker</blink>'
Markup(u'<strong>Hello &lt;blink&gt;hacker&lt;/blink&gt;!</strong>')
>>> Markup.escape('<blink>hacker</blink>')
Markup(u'&lt;blink&gt;hacker&lt;/blink&gt;')
>>> Markup('<em>Marked up</em> &raquo; HTML').striptags()
u'Marked up \xbb HTML'


\end{verbatim}


\end{frame}


\begin{frame}[t]{Rendering TemplatesÂ}

Changed in version 0.5: Autoescaping is no longer enabled for all templates.  The following
extensions for templates trigger autoescaping: .html, .htm,
.xml, .xhtml.  Templates loaded from a string will have
autoescaping disabled.

\end{frame}

\subsection{Accessing Request DataÂ}

\begin{frame}[t]{Accessing Request DataÂ}

For web applications it’s crucial to react to the data a client sends to
the server.  In Flask this information is provided by the global
request object.  If you have some experience with Python
you might be wondering how that object can be global and how Flask
manages to still be threadsafe.  The answer is context locals:

\end{frame}


\begin{frame}[t]{Accessing Request DataÂ}

Insider Information

\end{frame}


\begin{frame}[t]{Accessing Request DataÂ}

If you want to understand how that works and how you can implement
tests with context locals, read this section, otherwise just skip it.

\end{frame}


\begin{frame}[t]{Accessing Request DataÂ}

Certain objects in Flask are global objects, but not of the usual kind.
These objects are actually proxies to objects that are local to a specific
context.  What a mouthful.  But that is actually quite easy to understand.

\end{frame}


\begin{frame}[t]{Accessing Request DataÂ}

Imagine the context being the handling thread.  A request comes in and the
web server decides to spawn a new thread (or something else, the
underlying object is capable of dealing with concurrency systems other
than threads).  When Flask starts its internal request handling it
figures out that the current thread is the active context and binds the
current application and the WSGI environments to that context (thread).
It does that in an intelligent way so that one application can invoke another
application without breaking.

\end{frame}


\begin{frame}[t]{Accessing Request DataÂ}

So what does this mean to you?  Basically you can completely ignore that
this is the case unless you are doing something like unit testing.  You
will notice that code which depends on a request object will suddenly break
because there is no request object.  The solution is creating a request
object yourself and binding it to the context.  The easiest solution for
unit testing is to use the test\_request\_context()
context manager.  In combination with the with statement it will bind a
test request so that you can interact with it.  Here is an example:

\end{frame}


\begin{frame}[t,fragile]{Accessing Request DataÂ}


\begin{verbatim}

from flask import request

with app.test_request_context('/hello', method='POST'):
    # now you can do something with the request until the
    # end of the with block, such as basic assertions:
    assert request.path == '/hello'
    assert request.method == 'POST'


\end{verbatim}


\end{frame}


\begin{frame}[t]{Accessing Request DataÂ}

The other possibility is passing a whole WSGI environment to the
request\_context() method:

\end{frame}


\begin{frame}[t,fragile]{Accessing Request DataÂ}


\begin{verbatim}

from flask import request

with app.request_context(environ):
    assert request.method == 'POST'


\end{verbatim}


\end{frame}


\begin{frame}[t]{Accessing Request DataÂ}

The request object is documented in the API section and we will not cover
it here in detail (see request). Here is a broad overview of
some of the most common operations.  First of all you have to import it from
the flask module:

\end{frame}


\begin{frame}[t,fragile]{Accessing Request DataÂ}


\begin{verbatim}

from flask import request


\end{verbatim}


\end{frame}


\begin{frame}[t]{Accessing Request DataÂ}

The current request method is available by using the
method attribute.  To access form data (data
transmitted in a POST or PUT request) you can use the
form attribute.  Here is a full example of the two
attributes mentioned above:

\end{frame}


\begin{frame}[t,fragile]{Accessing Request DataÂ}


\begin{verbatim}

@app.route('/login', methods=['POST', 'GET'])
def login():
    error = None
    if request.method == 'POST':
        if valid_login(request.form['username'],
                       request.form['password']):
            return log_the_user_in(request.form['username'])
        else:
            error = 'Invalid username/password'
    # the code below is executed if the request method
    # was GET or the credentials were invalid
    return render_template('login.html', error=error)


\end{verbatim}


\end{frame}


\begin{frame}[t]{Accessing Request DataÂ}

What happens if the key does not exist in the form attribute?  In that
case a special KeyError is raised.  You can catch it like a
standard KeyError but if you don’t do that, a HTTP 400 Bad Request
error page is shown instead.  So for many situations you don’t have to
deal with that problem.

\end{frame}


\begin{frame}[t]{Accessing Request DataÂ}

To access parameters submitted in the URL (?key=value) you can use the
args attribute:

\end{frame}


\begin{frame}[t,fragile]{Accessing Request DataÂ}


\begin{verbatim}

searchword = request.args.get('key', '')


\end{verbatim}


\end{frame}


\begin{frame}[t]{Accessing Request DataÂ}

We recommend accessing URL parameters with get or by catching the
KeyError because users might change the URL and presenting them a 400
bad request page in that case is not user friendly.

\end{frame}


\begin{frame}[t]{Accessing Request DataÂ}

For a full list of methods and attributes of the request object, head over
to the request documentation.

\end{frame}


\begin{frame}[t]{Accessing Request DataÂ}

You can handle uploaded files with Flask easily.  Just make sure not to
forget to set the enctype="multipart/form-data" attribute on your HTML
form, otherwise the browser will not transmit your files at all.

\end{frame}


\begin{frame}[t]{Accessing Request DataÂ}

Uploaded files are stored in memory or at a temporary location on the
filesystem.  You can access those files by looking at the
files attribute on the request object.  Each
uploaded file is stored in that dictionary.  It behaves just like a
standard Python file object, but it also has a
save() method that allows you to store that
file on the filesystem of the server.  Here is a simple example showing how
that works:

\end{frame}


\begin{frame}[t,fragile]{Accessing Request DataÂ}


\begin{verbatim}

from flask import request

@app.route('/upload', methods=['GET', 'POST'])
def upload_file():
    if request.method == 'POST':
        f = request.files['the_file']
        f.save('/var/www/uploads/uploaded_file.txt')
    ...


\end{verbatim}


\end{frame}


\begin{frame}[t]{Accessing Request DataÂ}

If you want to know how the file was named on the client before it was
uploaded to your application, you can access the
filename attribute.  However please keep in
mind that this value can be forged so never ever trust that value.  If you
want to use the filename of the client to store the file on the server,
pass it through the secure\_filename() function that
Werkzeug provides for you:

\end{frame}


\begin{frame}[t,fragile]{Accessing Request DataÂ}


\begin{verbatim}

from flask import request
from werkzeug.utils import secure_filename

@app.route('/upload', methods=['GET', 'POST'])
def upload_file():
    if request.method == 'POST':
        f = request.files['the_file']
        f.save('/var/www/uploads/' + secure_filename(f.filename))
    ...


\end{verbatim}


\end{frame}


\begin{frame}[t]{Accessing Request DataÂ}

For some better examples, checkout the Uploading Files pattern.

\end{frame}


\begin{frame}[t]{Accessing Request DataÂ}

To access cookies you can use the cookies
attribute.  To set cookies you can use the
set\_cookie method of response objects.  The
cookies attribute of request objects is a
dictionary with all the cookies the client transmits.  If you want to use
sessions, do not use the cookies directly but instead use the
Sessions in Flask that add some security on top of cookies for you.

\end{frame}


\begin{frame}[t]{Accessing Request DataÂ}

Reading cookies:

\end{frame}


\begin{frame}[t,fragile]{Accessing Request DataÂ}


\begin{verbatim}

from flask import request

@app.route('/')
def index():
    username = request.cookies.get('username')
    # use cookies.get(key) instead of cookies[key] to not get a
    # KeyError if the cookie is missing.


\end{verbatim}


\end{frame}


\begin{frame}[t]{Accessing Request DataÂ}

Storing cookies:

\end{frame}


\begin{frame}[t,fragile]{Accessing Request DataÂ}


\begin{verbatim}

from flask import make_response

@app.route('/')
def index():
    resp = make_response(render_template(...))
    resp.set_cookie('username', 'the username')
    return resp


\end{verbatim}


\end{frame}


\begin{frame}[t]{Accessing Request DataÂ}

Note that cookies are set on response objects.  Since you normally
just return strings from the view functions Flask will convert them into
response objects for you.  If you explicitly want to do that you can use
the make\_response() function and then modify it.

\end{frame}


\begin{frame}[t]{Accessing Request DataÂ}

Sometimes you might want to set a cookie at a point where the response
object does not exist yet.  This is possible by utilizing the
Deferred Request Callbacks pattern.

\end{frame}


\begin{frame}[t]{Accessing Request DataÂ}

For this also see About Responses.

\end{frame}

\subsection{Redirects and ErrorsÂ}

\begin{frame}[t]{Redirects and ErrorsÂ}

To redirect a user to another endpoint, use the redirect()
function; to abort a request early with an error code, use the
abort() function:

\end{frame}


\begin{frame}[t,fragile]{Redirects and ErrorsÂ}


\begin{verbatim}

from flask import abort, redirect, url_for

@app.route('/')
def index():
    return redirect(url_for('login'))

@app.route('/login')
def login():
    abort(401)
    this_is_never_executed()


\end{verbatim}


\end{frame}


\begin{frame}[t]{Redirects and ErrorsÂ}

This is a rather pointless example because a user will be redirected from
the index to a page they cannot access (401 means access denied) but it
shows how that works.

\end{frame}


\begin{frame}[t]{Redirects and ErrorsÂ}

By default a black and white error page is shown for each error code.  If
you want to customize the error page, you can use the
errorhandler() decorator:

\end{frame}


\begin{frame}[t,fragile]{Redirects and ErrorsÂ}


\begin{verbatim}

from flask import render_template

@app.errorhandler(404)
def page_not_found(error):
    return render_template('page_not_found.html'), 404


\end{verbatim}


\end{frame}


\begin{frame}[t]{Redirects and ErrorsÂ}

Note the 404 after the render\_template() call.  This
tells Flask that the status code of that page should be 404 which means
not found.  By default 200 is assumed which translates to: all went well.

\end{frame}


\begin{frame}[t]{Redirects and ErrorsÂ}

See Error handlers for more details.

\end{frame}

\subsection{About ResponsesÂ}

\begin{frame}[t]{About ResponsesÂ}

The return value from a view function is automatically converted into a
response object for you.  If the return value is a string it’s converted
into a response object with the string as response body, a 200 OK
status code and a text/html mimetype.  The logic that Flask applies to
converting return values into response objects is as follows:

\end{frame}


\begin{frame}[t]{About ResponsesÂ}

If you want to get hold of the resulting response object inside the view
you can use the make\_response() function.

\end{frame}


\begin{frame}[t]{About ResponsesÂ}

Imagine you have a view like this:

\end{frame}


\begin{frame}[t,fragile]{About ResponsesÂ}


\begin{verbatim}

@app.errorhandler(404)
def not_found(error):
    return render_template('error.html'), 404


\end{verbatim}


\end{frame}


\begin{frame}[t]{About ResponsesÂ}

You just need to wrap the return expression with
make\_response() and get the response object to modify it, then
return it:

\end{frame}


\begin{frame}[t,fragile]{About ResponsesÂ}


\begin{verbatim}

@app.errorhandler(404)
def not_found(error):
    resp = make_response(render_template('error.html'), 404)
    resp.headers['X-Something'] = 'A value'
    return resp


\end{verbatim}


\end{frame}

\subsection{SessionsÂ}

\begin{frame}[t]{SessionsÂ}

In addition to the request object there is also a second object called
session which allows you to store information specific to a
user from one request to the next.  This is implemented on top of cookies
for you and signs the cookies cryptographically.  What this means is that
the user could look at the contents of your cookie but not modify it,
unless they know the secret key used for signing.

\end{frame}


\begin{frame}[t]{SessionsÂ}

In order to use sessions you have to set a secret key.  Here is how
sessions work:

\end{frame}


\begin{frame}[t,fragile]{SessionsÂ}


\begin{verbatim}

from flask import Flask, session, redirect, url_for, escape, request

app = Flask(__name__)

@app.route('/')
def index():
    if 'username' in session:
        return 'Logged in as %s' % escape(session['username'])
    return 'You are not logged in'

@app.route('/login', methods=['GET', 'POST'])
def login():
    if request.method == 'POST':
        session['username'] = request.form['username']
        return redirect(url_for('index'))
    return '''
        <form method="post">
            <p><input type=text name=username>
            <p><input type=submit value=Login>
        </form>
    '''

@app.route('/logout')
def logout():
    # remove the username from the session if it's there
    session.pop('username', None)
    return redirect(url_for('index'))

# set the secret key.  keep this really secret:
app.secret_key = 'A0Zr98j/3yX R~XHH!jmN]LWX/,?RT'


\end{verbatim}


\end{frame}


\begin{frame}[t]{SessionsÂ}

The escape() mentioned here does escaping for you if you are
not using the template engine (as in this example).

\end{frame}


\begin{frame}[t]{SessionsÂ}

How to generate good secret keys

\end{frame}


\begin{frame}[t]{SessionsÂ}

The problem with random is that it’s hard to judge what is truly random.  And
a secret key should be as random as possible.  Your operating system
has ways to generate pretty random stuff based on a cryptographic
random generator which can be used to get such a key:

\end{frame}


\begin{frame}[t,fragile]{SessionsÂ}


\begin{verbatim}

>>> import os
>>> os.urandom(24)
'\xfd{H\xe5<\x95\xf9\xe3\x96.5\xd1\x01O<!\xd5\xa2\xa0\x9fR"\xa1\xa8'

Just take that thing and copy/paste it into your code and you're done.


\end{verbatim}


\end{frame}


\begin{frame}[t]{SessionsÂ}

A note on cookie-based sessions: Flask will take the values you put into the
session object and serialize them into a cookie.  If you are finding some
values do not persist across requests, cookies are indeed enabled, and you are
not getting a clear error message, check the size of the cookie in your page
responses compared to the size supported by web browsers.

\end{frame}


\begin{frame}[t]{SessionsÂ}

Besides the default client-side based sessions, if you want to handle
sessions on the server-side instead, there are several
Flask extensions that support this.

\end{frame}

\subsection{Message FlashingÂ}

\begin{frame}[t]{Message FlashingÂ}

Good applications and user interfaces are all about feedback.  If the user
does not get enough feedback they will probably end up hating the
application.  Flask provides a really simple way to give feedback to a
user with the flashing system.  The flashing system basically makes it
possible to record a message at the end of a request and access it on the next
(and only the next) request.  This is usually combined with a layout
template to expose the message.

\end{frame}


\begin{frame}[t]{Message FlashingÂ}

To flash a message use the flash() method, to get hold of the
messages you can use get\_flashed\_messages() which is also
available in the templates.  Check out the Message Flashing
for a full example.

\end{frame}

\subsection{LoggingÂ}

\begin{frame}[t]{LoggingÂ}

New in version 0.3.

\end{frame}


\begin{frame}[t]{LoggingÂ}

Sometimes you might be in a situation where you deal with data that
should be correct, but actually is not.  For example you may have some client-side
code that sends an HTTP request to the server but it’s obviously
malformed.  This might be caused by a user tampering with the data, or the
client code failing.  Most of the time it’s okay to reply with 400 Bad
Request in that situation, but sometimes that won’t do and the code has
to continue working.

\end{frame}


\begin{frame}[t]{LoggingÂ}

You may still want to log that something fishy happened.  This is where
loggers come in handy.  As of Flask 0.3 a logger is preconfigured for you
to use.

\end{frame}


\begin{frame}[t]{LoggingÂ}

Here are some example log calls:

\end{frame}


\begin{frame}[t,fragile]{LoggingÂ}


\begin{verbatim}

app.logger.debug('A value for debugging')
app.logger.warning('A warning occurred (%d apples)', 42)
app.logger.error('An error occurred')


\end{verbatim}


\end{frame}


\begin{frame}[t]{LoggingÂ}

The attached logger is a standard logging
Logger, so head over to the official logging
documentation for more
information.

\end{frame}


\begin{frame}[t]{LoggingÂ}

Read more on Application Errors.

\end{frame}

\subsection{Hooking in WSGI MiddlewaresÂ}

\begin{frame}[t]{Hooking in WSGI MiddlewaresÂ}

If you want to add a WSGI middleware to your application you can wrap the
internal WSGI application.  For example if you want to use one of the
middlewares from the Werkzeug package to work around bugs in lighttpd, you
can do it like this:

\end{frame}


\begin{frame}[t,fragile]{Hooking in WSGI MiddlewaresÂ}


\begin{verbatim}

from werkzeug.contrib.fixers import LighttpdCGIRootFix
app.wsgi_app = LighttpdCGIRootFix(app.wsgi_app)


\end{verbatim}


\end{frame}

\subsection{Using Flask ExtensionsÂ}

\begin{frame}[t]{Using Flask ExtensionsÂ}

Extensions are packages that help you accomplish common tasks. For
example, Flask-SQLAlchemy provides SQLAlchemy support that makes it simple
and easy to use with Flask.

\end{frame}


\begin{frame}[t]{Using Flask ExtensionsÂ}

For more on Flask extensions, have a look at Flask Extensions.

\end{frame}

\subsection{Deploying to a Web ServerÂ}

\begin{frame}[t]{Deploying to a Web ServerÂ}

Ready to deploy your new Flask app? Go to Deployment Options.

\end{frame}


\end{document}

